\documentclass[11pt,a4paper]{article}
\usepackage{fontspec,amsmath,amssymb,amsthm,microtype}
\setmainfont{lmroman10-regular.otf}[BoldFont=lmroman10-bold.otf,ItalicFont=lmroman10-italic.otf,BoldItalicFont=lmroman10-bolditalic.otf]
\setmonofont{lmmono10-regular.otf}[Scale=MatchLowercase]
\usepackage[margin=25mm]{geometry}
\usepackage[unicode,hidelinks]{hyperref}
\setlength{\emergencystretch}{3em}
\setlength{\parskip}{0.4em}
\setlength{\parindent}{0pt}
\allowdisplaybreaks[2]
\raggedbottom
\title{\texorpdfstring{Absorb the Hessian cutoff and accumulate the second-\allowbreak{}order error}{Absorb the Hessian cutoff and accumulate the second-order error}}
\author{LangevinLean mathematical exposition}
\date{}
\begin{document}\maketitle
\noindent Source module: \texttt{Kinetic\allowbreak{}Langevin/\allowbreak{}Second\allowbreak{}Order\allowbreak{}Recurrence.\allowbreak{}lean}. The accompanying HTML pairs this mathematical exposition with the exact Lean statements and records the scope of the formalization.


The repaired local comparison has a random-Hessian term proportional to \(M_2\sqrt\alpha h^2\). Choosing the cutoff \(\alpha=(m/(4M_2h))^2\) absorbs it into the contraction gap. This module proves the choice, both Gaussian tails, the affine accumulation, the position projection, and the exact dimension range for the manuscript’s final simplification.


\section{\texorpdfstring{Choose the cutoff and retain both tail estimates}{Choose the cutoff and retain both tail estimates}}


For \(m,M_2,h>0\), the positive square root of \(\alpha=(m/(4M_2h))^2\) is \(m/(4M_2h)\). Thus \(M_2\sqrt\alpha h=m/4\). The source dimension restriction \(4\sqrt{5p}M_2h\le m\) implies \(\alpha\ge5p\). It follows that \((\alpha-p)/8\ge p/2\), giving the printed tail. Also \(p\le\alpha/5\) implies \((\alpha-p)/8\ge\alpha/10=m^2/(160M_2^2h^2)\), giving the stronger source footnote. Lean proves the square-root identity, multiplies through positive denominators, compares nonnegative squares, and applies monotonicity of the exponential.



{\raggedright\noindent Lean declarations: \texttt{second\allowbreak{}Order\allowbreak{}Cutoff}, \texttt{second\allowbreak{}Order\allowbreak{}Cutoff\_\allowbreak{}spec}, \texttt{second\allowbreak{}Order\allowbreak{}Cutoff\_\allowbreak{}tail}.\par}

\section{\texorpdfstring{Sum the actual one-\allowbreak{}step comparison and project to position}{Sum the actual one-step comparison and project to position}}


The local bound has coefficient \(1-mh/(2\gamma)+M_2\sqrt\alpha h^2/\gamma\). The cutoff makes this at most \(q=1-mh/(4\gamma)\). The additive local terms are \(\sqrt2h^2B e^{-(\alpha-p)/8}+h^3C/4\), where \(B=M-m\) and \(C=M_2\sqrt{p^2+2p}+M\sqrt{Mp}\). The checked affine recurrence divides this budget by \(1-q=mh/(4\gamma)\). Position projection multiplies by \(\sqrt2/\gamma\), and equal initial velocities give \(e_0\le\gamma d_0\). The exact resulting bias is \(\sqrt2h^2C/m+(8B/m)h e^{-(\alpha-p)/8}\), with the desired transient \(\sqrt2q^nd_0\). The recurrence is a reusable scalar lemma; the actual chain endpoint supplies its one-step inequality separately.



{\raggedright\noindent Lean declarations: \texttt{second\allowbreak{}Order\_\allowbreak{}coupling\_\allowbreak{}recurrence\_\allowbreak{}bound}, \texttt{second\allowbreak{}Order\_\allowbreak{}projected\_\allowbreak{}bound}.\par}

\section{\texorpdfstring{State the dimension restriction explicitly}{State the dimension restriction explicitly}}


For \(p\ge2\), \(p^2+2p\le2p^2\), hence \(\sqrt2\sqrt{p^2+2p}\le2p\). Distributing this through \(C\) gives the printed Hessian term \(2M_2p\), while \(\sqrt2\sqrt{Mp}=\sqrt{2Mp}\) gives the score term. The supplied proof’s simplification does not justify the same expression at \(p=1\). The general theorem therefore retains the exact fourth-moment expression there; the printed simplification is an explicit \(p\ge2\) corollary. Lean compares nonnegative squares and never uses the invalid inequality in dimension one.



{\raggedright\noindent Lean declarations: \texttt{second\allowbreak{}Order\_\allowbreak{}source\_\allowbreak{}dimension\_\allowbreak{}bound}.\par}

\end{document}
